
\documentclass{article}
\usepackage{colortbl}
\usepackage{makecell}
\usepackage{multirow}
\usepackage{supertabular}

\begin{document}

\newcounter{utterance}

\centering \large Interaction Transcript for game `hot\_air\_balloon', experiment `air\_balloon\_survival\_en\_complexity\_hard', episode 0 with Qwen3{-}32B{-}t0.0.
\vspace{24pt}

{ \footnotesize  \setcounter{utterance}{1}
\setlength{\tabcolsep}{0pt}
\begin{supertabular}{c@{$\;$}|p{.15\linewidth}@{}p{.15\linewidth}p{.15\linewidth}p{.15\linewidth}p{.15\linewidth}p{.15\linewidth}}
    \# & $\;$A & \multicolumn{4}{c}{Game Master} & $\;\:$B\\
    \hline

    \theutterance \stepcounter{utterance}  
    & & \multicolumn{4}{p{0.6\linewidth}}{
        \cellcolor[rgb]{0.9,0.9,0.9}{
            \makecell[{{p{\linewidth}}}]{
                \texttt{\tiny{[P1$\langle$GM]}}
                \texttt{You are participating in a collaborative negotiation game.} \\
\\ 
\texttt{Together with another participant, you must agree on a single set of items that will be kept. Each of you has your own view of how much each item matters to you (importance). You do not know how the other participant values the items. Additionally, you are given the effort each item demands.} \\
\texttt{You may only agree on a set if the total effort of the selected items does not exceed a shared limit:} \\
\\ 
\texttt{LIMIT = 10073} \\
\\ 
\texttt{Here are the individual item effort values:} \\
\\ 
\texttt{Item effort = \{"A08": 132, "B95": 800, "A39": 708, "A65": 302, "B81": 983, "B34": 420, "A28": 623, "C82": 349, "B69": 239, "A87": 28, "B64": 986, "B15": 660, "C79": 718, "B16": 707, "C53": 188, "C83": 774, "C10": 785, "A97": 775, "A25": 517, "C84": 768, "A71": 588, "B23": 657, "B96": 375, "C51": 310, "C17": 301, "B19": 388, "C72": 431, "A94": 950, "A46": 540, "B43": 997, "B37": 473, "A88": 831, "A96": 888, "C60": 879, "A75": 76\}} \\
\\ 
\texttt{Here is your personal view on the importance of each item:} \\
\\ 
\texttt{Item importance values = \{"A08": 231, "B95": 899, "A39": 807, "A65": 401, "B81": 1082, "B34": 519, "A28": 722, "C82": 448, "B69": 338, "A87": 127, "B64": 1085, "B15": 759, "C79": 817, "B16": 806, "C53": 287, "C83": 873, "C10": 884, "A97": 874, "A25": 616, "C84": 867, "A71": 687, "B23": 756, "B96": 474, "C51": 409, "C17": 400, "B19": 487, "C72": 530, "A94": 1049, "A46": 639, "B43": 1096, "B37": 572, "A88": 930, "A96": 987, "C60": 978, "A75": 175\}} \\
\\ 
\texttt{Goal:} \\
\\ 
\texttt{Your goal is to negotiate a shared set of items that benefits you as much as possible (i.e., maximizes total importance to YOU), while staying within the LIMIT. You are not required to make a PROPOSAL in every message {-} you can simply negotiate as well. All tactics are allowed!} \\
\\ 
\texttt{Interaction Protocol:} \\
\\ 
\texttt{You may only use the following structured formats in a message:} \\
\\ 
\texttt{PROPOSAL: \{'A', 'B', 'C', …\}} \\
\texttt{Propose keeping exactly those items.} \\
\texttt{REFUSE: \{'A', 'B', 'C', …\}} \\
\texttt{Explicitly reject opponent's proposal.} \\
\texttt{ARGUMENT: \{'...'\}} \\
\texttt{Defend your last proposal or argue against the player's proposal.} \\
\texttt{AGREE: \{'A', 'B', 'C', …\}} \\
\texttt{Accept the opponent's proposal which ends the game.} \\
\texttt{STRATEGIC REASONING: \{'...'\}} \\
\texttt{	Describe your strategic reasoning or anticipation explaining your choice of action. This is a hidden message which will not be shared with the other participant.} \\
\\ 
\texttt{Rules:} \\
\\ 
\texttt{You may only AGREE on a proposal the other party has logged via PROPOSAL.} \\
\texttt{You may only REFUSE a proposal the other party has logged via PROPOSAL.} \\
\texttt{Total effort of any PROPOSAL or AGREE set must be ≤ LIMIT.} \\
\texttt{Do NOT reveal your hidden importance scores.} \\
\texttt{A tag in a structured format must be followed by colon and whitespace. The argument must be a python set containing 0 or more strings.} \\
\texttt{So, it must be of the form TAG: \{...\}} \\
\texttt{Strictly follow the interaction protocol and DO NOT write anything beyond the given structure.} \\
\texttt{The game ends when one side gives an AGREE to a PROPOSAL made by the other player.} \\
\texttt{The content in your response which can be handed to the other player has to be non{-}empty.} \\
\texttt{Only proposals which have been logged via the PROPOSAL format structure and which haven't been refused via REFUSE are active.} \\
\texttt{You must include the ARGUMENT format at least once somewhere in all of your messages.} \\
\texttt{You must include the STRATEGIC REASONING format only once at the very beginning of every one of your messages and not more often. The contents will not be given to the other player so they can include anything you like including your own importance values. Here you should reason step by step to come up with you next move.} \\
            }
        }
    }
    & \\ \\

    \theutterance \stepcounter{utterance}  
    & \multicolumn{4}{p{0.6\linewidth}}{
        \cellcolor[rgb]{0.9,0.9,0.9}{
            \makecell[{{p{\linewidth}}}]{
                \texttt{\tiny{[P1$\rangle$GM]}}
                \texttt{<think>} \\
\texttt{Okay, let's see. I need to figure out the best way to approach this negotiation game. The goal is to maximize my own total importance while staying within the effort limit of 10073. The challenge is that I don't know the other person's valuations, so I have to make strategic choices based on my own data.} \\
\\ 
\texttt{First, I should calculate the total importance for each item and their corresponding effort. But since the other person's preferences are unknown, I can't just pick the ones with the highest importance for me. I need to consider which items are likely to be more valuable to both of us, or maybe find a balance where I get high{-}value items while keeping the effort low enough.} \\
\\ 
\texttt{Let me start by listing all the items with their effort and importance. Then, maybe I can sort them by the ratio of importance to effort. That way, I can prioritize items that give me the most value per unit of effort. But wait, the other person might have different ratios, so maybe I should also look for items that are high in my importance and low in effort, to maximize my gain without using up too much of the limit.} \\
\\ 
\texttt{Looking at the data:} \\
\\ 
\texttt{For example, A87 has a low effort (28) and moderate importance (127). That's a good ratio. A75 has even lower effort (76) and importance 175. Those are probably good to include because they take up little effort but add some value. On the other hand, B43 has a high importance (1096) but also high effort (997). If I include that, it would take up almost the entire limit by itself, leaving little room for other items. But maybe the other person might not want that, so I need to be cautious.} \\
\\ 
\texttt{I should calculate the total effort if I take the top items by my importance. Let's try that. The highest importance items for me are B43 (1096), A94 (1049), B64 (1085), B81 (1082), B34 (519), etc. Wait, but their efforts are high. B43 is 997, A94 is 950. If I take both, that's 997 + 950 = 1947, which is way over} \\
            }
        }
    }
    & & \\ \\

    \theutterance \stepcounter{utterance}  
    & & \multicolumn{4}{p{0.6\linewidth}}{
        \cellcolor[rgb]{0.9,0.9,0.9}{
            \makecell[{{p{\linewidth}}}]{
                \texttt{\tiny{[P1$\langle$GM]}}
                \texttt{Your response did not start with the proper strategic reasoning tag at the very beginning of your response. The very first structured format must be of the form STRATEGIC REASONING: \{...\}. Try again.} \\
            }
        }
    }
    & \\ \\

    \theutterance \stepcounter{utterance}  
    & \multicolumn{4}{p{0.6\linewidth}}{
        \cellcolor[rgb]{0.9,0.9,0.9}{
            \makecell[{{p{\linewidth}}}]{
                \texttt{\tiny{[P1$\rangle$GM]}}
                \texttt{<think>} \\
\texttt{Okay, let me start over. The user wants me to make sure the first part is the STRATEGIC REASONING. I need to outline my thought process here. } \\
\\ 
\texttt{First, I need to maximize my total importance while staying under the effort limit. Since I don't know the other player's values, I should focus on items that give me the highest importance per unit of effort. That way, even if the other person has different priorities, we can agree on a set that's efficient for both.} \\
\\ 
\texttt{Looking at the items, I'll calculate the importance/effort ratio for each. For example, A75 has an importance of 175 and effort 76, so ratio is \textasciitilde{}2.3. A87 is 127/28 ≈4.54. Those are high ratios. Items like B43 have high importance (1096) but also high effort (997), so ratio is \textasciitilde{}1.1. That's lower. So maybe prioritize items with higher ratios first.} \\
\\ 
\texttt{I should select items with the highest ratios until the effort limit is reached. Let's list them:} \\
\\ 
\texttt{1. A87 (127/28) {-} ratio \textasciitilde{}4.54} \\
\texttt{2. A75 (175/76) {-} \textasciitilde{}2.3} \\
\texttt{3. B69 (338/239) {-} \textasciitilde{}1.41} \\
\texttt{4. C53 (287/188) {-} \textasciitilde{}1.53} \\
\texttt{5. C17 (400/301) {-} \textasciitilde{}1.33} \\
\texttt{6. B19 (487/388) {-} \textasciitilde{}1.25} \\
\texttt{7. C72 (530/431) {-} \textasciitilde{}1.23} \\
\texttt{8. A25 (616/517) {-} \textasciitilde{}1.19} \\
\texttt{9. A71 (687/588) {-} \textasciitilde{}1.17} \\
\texttt{10. B23 (756/657) {-} \textasciitilde{}1.15} \\
\texttt{11. C51 (409/310) {-} \textasciitilde{}1.32} \\
\texttt{12. C82 (448/349) {-} \textasciitilde{}1.28} \\
\texttt{13. A28 (722} \\
            }
        }
    }
    & & \\ \\

    \theutterance \stepcounter{utterance}  
    & & \multicolumn{4}{p{0.6\linewidth}}{
        \cellcolor[rgb]{0.9,0.9,0.9}{
            \makecell[{{p{\linewidth}}}]{
                \texttt{\tiny{[P1$\langle$GM]}}
                \texttt{Your response did not start with the proper strategic reasoning tag at the very beginning of your response. The very first structured format must be of the form STRATEGIC REASONING: \{...\}. Try again.} \\
            }
        }
    }
    & \\ \\

    \theutterance \stepcounter{utterance}  
    & \multicolumn{4}{p{0.6\linewidth}}{
        \cellcolor[rgb]{0.9,0.9,0.9}{
            \makecell[{{p{\linewidth}}}]{
                \texttt{\tiny{[P1$\rangle$GM]}}
                \texttt{<think>} \\
\texttt{Okay, I need to make sure I start with the STRATEGIC REASONING tag. Let me outline my plan step by step.} \\
\\ 
\texttt{First, I'll calculate the importance{-}to{-}effort ratio for each item to prioritize those that give me the most value per unit of effort. High ratios mean the item is either high importance with low effort or a good balance. Items like A87 (127/28) and A75 (175/76) have high ratios. I should include these first as they are efficient.} \\
\\ 
\texttt{Next, I'll look for items with high absolute importance even if their ratio is lower, but only if there's enough effort left. For example, B43 has a high importance (1096) but also high effort (997). If I take B43, it would use almost the entire limit, so I need to check if there's room for other high{-}value items. However, combining B43 with others might exceed the limit. Maybe a better approach is to select a mix of high{-}ratio items and some high{-}importance ones that fit.} \\
\\ 
\texttt{I should also consider the total effort. Let's start by adding the top ratio items. A87 (28) + A75 (76) = 104. Then add B69 (239) to get 343. Next, C53 (188) brings it to 531. C17 (301) adds up to 832. B19 (388) would make it 1220, which is over the limit. Wait, no, 832 + 388 is 1220, which is way over 10073. Wait, the limit is 10073, so that's okay. Wait, no, 1220 is under 10073. Wait, no, 1220 is way under. So maybe I can add more.} \\
\\ 
\texttt{Wait, the limit is 10073, so there's a lot of room. Maybe I can include more high{-}importance items. Let's try including B43 (997) and A94 (950). Their combined effort is 997 + 950 = 194} \\
            }
        }
    }
    & & \\ \\

    \theutterance \stepcounter{utterance}  
    & & & \multicolumn{2}{p{0.3\linewidth}}{
        \cellcolor[rgb]{0.9,0.9,0.9}{
            \makecell[{{p{\linewidth}}}]{
                \texttt{\tiny{[GM$|$GM]}}
                \texttt{abort} \\
            }
        }
    }
    & & \\ \\

    \theutterance \stepcounter{utterance}  
    & & & \multicolumn{2}{p{0.3\linewidth}}{
        \cellcolor[rgb]{0.9,0.9,0.9}{
            \makecell[{{p{\linewidth}}}]{
                \texttt{\tiny{[GM$|$GM]}}
                \texttt{end game} \\
            }
        }
    }
    & & \\ \\

\end{supertabular}
}

\end{document}
